% ---------------------------------------------------------------------------
% Comparison with other methods (accuracy and computational cost).
% Table rows in data/tab_methods.tex are written by make_guide_figures.m.
% ---------------------------------------------------------------------------
\subsection{Comparison with other methods: accuracy and computational cost}
\label{sec:comparison}

To place the solver in context, the same Sod problem was solved on the same grids with
five alternative methods:
\begin{itemize}
  \item the \textbf{same DBM} with three other reconstructions: first-order upwind (no
        slope), MUSCL with minmod, and MUSCL with MC; and
  \item a \textbf{conventional finite-volume Euler solver}, which evolves density,
        momentum and energy directly, using the HLLC approximate Riemann
        solver~\citep{ToroSpruceSpeares1994,Toro2009} with either first-order reconstruction
        or the same generalized-minmod MUSCL reconstruction ($\theta = 1.2$) and SSP-RK3
        time stepping at a standard CFL number of 0.5.
\end{itemize}
All computations used the same computer, one processor core and MATLAB R2024b.  The Euler
times are the best of three measured runs.  The DBM times are the measured cost per cell per
time step ($4.05\times10^{-7}$\,s for MUSCL and $2.74\times10^{-7}$\,s for first order,
obtained from isolated runs on the two coarsest grids) multiplied by the number of cells and
time steps; this cost was verified to be the same on both grids to within 0.4\%.

\begin{table}[H]
\centering
\caption{Accuracy and cost of the six methods.  Orders are least-squares fits over the grids
3{,}125--25{,}000 (common to all methods); errors and cost are for 25{,}000 cells.}
\label{tab:methods}
\small
\begin{tabular}{@{}lcccccc@{}}
\toprule
 & \multicolumn{2}{c}{Observed order ($\rho$)} & \multicolumn{2}{c}{Error at 25{,}000 cells} & \multicolumn{2}{c}{Cost at 25{,}000 cells} \\
\cmidrule(lr){2-3}\cmidrule(lr){4-5}\cmidrule(l){6-7}
Method & $L_1$ & $L_2$ & $L_1(\rho)$ & $L_2(\rho)$ & time steps & time [s] \\
\midrule
\tabinput{data/tab_methods.tex}
\bottomrule
\end{tabular}
\end{table}

\begin{figure}[H]
\centering
\includegraphics[width=\textwidth]{G09_order_comparison.pdf}
\caption{Density error versus cell size for the six methods: $L_1$ (a) and $L_2$ (b).}
\label{fig:ordercomp}
\end{figure}

\begin{figexplain}{\Cref{fig:ordercomp}}
\begin{fx}
\item[What is shown.] The density error of every method on every grid.  Blue curves are
  DBM variants (the dark solid curve is the solver studied here); red and orange curves are
  the conventional Euler solver.  Hollow markers denote first-order methods.
\item[How to read it.] Lower curves are more accurate at a given grid; steeper curves improve
  faster with refinement.  The dotted lines give reference slopes of 1 and $1/2$.
\item[What it means.] The methods separate into two groups.  The two first-order methods
  (hollow markers) have errors three to seven times larger, with the gap widening under
  refinement, and converge more slowly ($L_1$ order 0.64).  All second-order methods cluster together with $L_1$ orders of
  0.80--0.92, whether they are kinetic (DBM) or conventional (Euler).  At equal resolution
  the density error of the present DBM is about one third larger than that of the
  conventional MUSCL-HLLC solver.  The order of accuracy is therefore determined by the
  reconstruction, not by whether the gas is modelled kinetically or with the Euler equations.
\end{fx}
\end{figexplain}

\begin{figure}[H]
\centering
\includegraphics[width=\textwidth]{G10_accuracy_vs_cost.pdf}
\caption{(a) Density error versus computing time for the six methods.  (b) Number of time
steps required by the DBM and the Euler solver on each grid.}
\label{fig:cost}
\end{figure}

\begin{figexplain}{\Cref{fig:cost}}
\begin{fx}
\item[What is shown.] Panel (a) is a work--precision diagram: each point is one grid, placed by
  its computing time (horizontal) and its density error (vertical).  Panel (b) shows how many
  time steps each solver needed on each grid.
\item[How to read it.] In (a), the most efficient method is the one whose curve lies furthest to
  the lower left, meaning the lowest error for a given time.  A horizontal line through the
  plot compares the time each method needs to reach the same error.
\item[What it means.] The conventional Euler solvers are two to three orders of magnitude
  faster than the DBM for the same accuracy.  Panel (b) explains why.  The cost per cell per
  time step is similar for both solvers ($4.1\times10^{-7}$\,s for the DBM versus about
  $2.9\times10^{-7}$\,s for MUSCL-HLLC), but the DBM needs 120{,}001 steps on every grid
  because its time step is fixed by the collision limit $\Delta t = \tau/4$, whereas the Euler
  solver's step is set by the wave speed and needs only 102 steps on the coarsest grid and
  1{,}643 on the finest.  The gap therefore narrows with refinement, from about 780 times at
  6{,}250 cells to about 100 times at 50{,}000 cells.
\end{fx}
\end{figexplain}

\paragraph{Quantitative summary.}
\begin{itemize}
  \item \textbf{Second-order versus first-order DBM.}  At the same grid, MUSCL with
        $\theta = 1.2$ reduces the density error by a factor of 5.4 for only 1.5 times the
        cost per step.  Reaching the same accuracy with first-order upwinding would require
        roughly 14 times as many cells and about 9 times the computing time.
  \item \textbf{Choice of limiter.}  MC is about 15\% more accurate on average than
        $\theta = 1.2$ at identical cost, but produces the temperature overshoot described in
        \cref{sec:limiters}; minmod has about 35\% larger error.
  \item \textbf{DBM versus conventional Euler solver.}  On 50{,}000 cells the DBM reaches
        $L_1(\rho) = 2.5\times10^{-5}$ in an estimated 2{,}430\,s.  The MUSCL-HLLC solver reaches the
        same error on about 32{,}000 cells in roughly 10\,s, about 250 times less.  At equal
        resolution its density error is 25\% smaller (at 25{,}000 cells).
  \item \textbf{Source of the cost.}  The difference is entirely due to the number of time
        steps.  If the DBM time step were set by the transport CFL condition (CFL $= 0.5$)
        instead of by $\tau/4$, it would need about 2{,}250 steps on 50{,}000 cells instead of
        120{,}001, a reduction by a factor of about 50.  This requires treating the collision
        term implicitly~\citep{PareschiRusso2005} or using a scheme that couples transport and
        collision~\citep{Guo2013}.
\end{itemize}

\paragraph{Context from the literature.}
\Cref{tab:literature} places these results alongside other widely used shock-capturing
approaches.  An important reference point is the study of Greenough and
Rider~\citep{Greenough2004}, who compared a fifth-order WENO scheme~\citep{JiangShu1996}
with a second-order piecewise-linear Godunov (MUSCL-type) scheme.  For problems with
discontinuities, including the Sod problem, both schemes converged at approximately first
order; on Sod the second-order scheme had about half the error of WENO and was about six
times cheaper at equal resolution.  Comparative studies of many schemes on standard Euler
test problems~\citep{LiskaWendroff2003} similarly show that no method avoids the
accuracy limit imposed by discontinuities.  A second-order limited reconstruction is
therefore a sound choice for shock-dominated flows, and the present DBM uses one.  Its
efficiency disadvantage arises from the explicit treatment of collisions, not from the
spatial scheme.

\begin{table}[H]
\centering
\caption{Qualitative comparison of shock-capturing approaches for problems such as Sod.}
\label{tab:literature}
\small
\begin{tabularx}{\textwidth}{@{}p{3.3cm}cX@{}}
\toprule
Method & Formal order & Accuracy and cost on shock problems \\
\midrule
First-order upwind / Godunov & 1 & Robust but strongly diffusive; errors 3--7 times larger than
  second-order methods at equal resolution (\cref{tab:methods}). \\
MUSCL / TVD with limiter~\citep{vanLeer1979,Sweby1984} & 2 & Standard choice; convergence limited
  by discontinuities to about first order in $L_1$; cheapest route to a given accuracy on
  shock problems~\citep{Greenough2004}. \\
WENO5~\citep{JiangShu1996} & 5 & Far more accurate for smooth flow; on Sod, about first-order
  convergence with larger error and about six times higher cost than a second-order Godunov
  scheme at equal resolution~\citep{Greenough2004}. \\
Conventional Euler MUSCL-HLLC (this study) & 2 & $L_1(\rho)$ order 0.92; about 250 times cheaper
  than the DBM for equal accuracy on this problem. \\
FVDBM, explicit collision (this study) & 2 & Accuracy close to the conventional second-order
  solver; time step limited to $\tau/4$, giving 120{,}001 steps on every grid.  Provides kinetic
  information not available from the Euler equations. \\
Kinetic schemes with implicit or coupled collision~\citep{PareschiRusso2005,Guo2013} & 2 &
  Remove the $\Delta t \propto \tau$ restriction; the expected route to making the DBM
  competitive in cost. \\
\bottomrule
\end{tabularx}
\end{table}
